\section{DivMod Solution Invariance Under Linear Shift}
\label{sec:shift-invariance}

Moving one copy of the divisor between quotient and remainder preserves the
represented dividend; normalization therefore reaches the same final state.

\begin{equation*}
\begin{aligned}
\forall a,b,q,r \in \mathbb{Z},\; b \neq 0,\; a &= bq + r \\
a &= b(q+1) + (r-b) \\
a &= b(q-1) + (r+b) \\
\operatorname{DivMod}(a,b,q+1,r-b).\operatorname{solve}
  &= \operatorname{DivMod}(a,b,q,r).\operatorname{solve} \\
\operatorname{DivMod}(a,b,q-1,r+b).\operatorname{solve}
  &= \operatorname{DivMod}(a,b,q,r).\operatorname{solve}
\end{aligned}
\end{equation*}

This invariant is verified for the
\href{https://github.com/thiagomata/prime-numbers/blob/modulo-article-v1.0.0/src/main/scala/v1/chapter2/div/properties/ModIdempotence.scala#assertDivModWithMoreDivAndLessModSameSolution}{positive shift}
and
\href{https://github.com/thiagomata/prime-numbers/blob/modulo-article-v1.0.0/src/main/scala/v1/chapter2/div/properties/ModIdempotence.scala#assertDivModWithLessDivAndMoreModSameSolution}{negative shift}.

\subsection{Creating the Division and Modulo Operations}
\label{sec:creating-operations}

Using the normalized \texttt{DivMod} value,
\href{https://github.com/thiagomata/prime-numbers/blob/modulo-article-v1.0.0/src/main/scala/v1/chapter2/div/Calc.scala}{\texttt{Calc.scala}}
defines $\operatorname{div}$ and $\operatorname{mod}$ as the quotient and
remainder projections of the solved state. Starting from
$\operatorname{DivMod}(a,b,0,a)$, let:

\begin{equation*}
\begin{aligned}
S &:= \operatorname{DivMod}(a,b,0,a).\operatorname{solve} \\
q &:= S.\operatorname{div} \\
r &:= S.\operatorname{mod} \\
\operatorname{div}(a,b) &:= q \\
\operatorname{mod}(a,b) &:= r
\end{aligned}
\end{equation*}

So $\operatorname{div}(a,b)$ names the normalized quotient, while
$\operatorname{mod}(a,b)$ names the normalized remainder. In the source code
these are the \texttt{div} and \texttt{mod} fields of the solved
\texttt{DivMod}; in the article notation, $q$ and $r$ keep the quotient and
remainder roles separate from the operation names.

The article uses both functional and infix notation for the same operations:
$\operatorname{div}(x,y)$ and $x \operatorname{div} y$ are equivalent, as are
$\operatorname{mod}(x,y)$ and $x \operatorname{mod} y$. Functional notation is
useful when the operation is nested inside another expression; infix notation
keeps simple algebraic identities closer to the traditional presentation.
